The HG002 example separates annotation by how often it changes rather than
by where it comes from. Core annotation follows Ensembl releases and plugin
data follow their providers; both are compiled once into indexed Parquet
caches, so a weekly ClinVar release rebuilds a 0.1\,GiB plugin cache, not
the core cache. Data that a particular question needs---gene constraint,
curated gene--disease relationships, genome stratifications or the genotypes
of relatives---need no cache at all: they are read from their URLs when the
query runs and combined with the annotation through joins and interval
operations. In a VEP-based workflow, each such source is configured as a
custom annotation or plugin before annotation, or processed afterwards from
an annotated VCF with separate tools. Keeping the third tier outside the
annotation engine lets analyses change without changing, or re-running,
the annotation.
